\section{Information retained by a score label}

For a fixed score law and deterministic release, the sharp interval in \cref{obj:def:sharp-ate-set} can vary with the released outcome law. The worst-case ambiguity \leanref{sym:D}{\ensuremath{D_H(g)}}, the largest such interval width over the causal laws in \cref{obj:def:causal-law-class}, measures the information retained by the label across outcome laws.

Three criteria answer different questions. The sharp width measures identification for a specified released law; worst-case ambiguity maximizes that width over released laws at a fixed score law and release. With trial sampling and a known score law, \cref{obj:thm:minimax-honest-length} studies expected honest interval length. With an independently sampled score law, \cref{obj:def:external-excess-risk} studies expected length in excess of the population sharp interval. The latter two criteria include sampling uncertainty; the sharp width does not.

\begin{definitionv}{def:worst-case-ambiguity}[Worst-case ambiguity $D_H(g)$]
\[
D_H(g)
=
\sup_{P\in\mathfrak P(H,g)}
\operatorname{len}\!\left\{
I\bigl(\mathcal L_P(R,A,Y);H,g\bigr)
\right\}.
\]
\end{definitionv}

The supremum has a cellwise expression in terms of the score law and release rule. The next result also gives a released outcome law at which the width is attained.

\begin{theoremv}{thm:all-label-ambiguity}[Exact all-label ambiguity width]
Suppose that:
\begin{itemize}
\item \textbf{(Overlap.)} The overlap condition in \cref{obj:ass:overlap} holds.
\item \textbf{(Score law.)} \(H\) is a probability law on \(\mathcal E=[\varepsilon,1-\varepsilon]\).
\item \textbf{(Release.)} \(g:\mathcal E\to\mathcal R_J\) is measurable.
\end{itemize}
Write \(C_r=\{e:g(e)=r\}\), \(p_1(e)=e\), \(p_0(e)=1-e\), and \(q_{ar}=\int_{C_r}p_a(e)\,H(de)\). Let \(P_{\mathrm{rel}}\) be the released law
\[
P_{\mathrm{rel}}
=\sum_{r\in\mathcal R_J}\sum_{a\in\{0,1\}}
\frac{q_{ar}}{2}\bigl(\delta_{(r,a,0)}+\delta_{(r,a,1)}\bigr).
\]
Thus its observed outcome law \(F_{ar}\) is Bernoulli\((1/2)\) in every positive-mass arm-label cell.

The worst-case ambiguity \(D_H(g)\) of \cref{obj:def:worst-case-ambiguity} satisfies
\[
D_H(g)=
\sum_{r\in\mathcal R_J}
\left[
\inf_{c\in[1/(1-\varepsilon),\,1/\varepsilon]}
\int_{C_r}|1-ce|\,H(de)
+
\inf_{c\in[1/(1-\varepsilon),\,1/\varepsilon]}
\int_{C_r}|1-c(1-e)|\,H(de)
\right].
\]

For each \(a\in\{0,1\}\) and each cell with \(q_{ar}>0\), let \(E_{ar}\) have law \(p_a(e)\mathbf 1_{C_r}(e)H(de)/q_{ar}\), and let \(W_{ar}=1/p_a(E_{ar})\). Define
\[
\underline\mu_a
=\sum_{r:q_{ar}>0}q_{ar}\int_0^1
Q_{F_{ar}}(u)Q_{W_{ar}}(1-u)\,du,
\qquad
\overline\mu_a
=\sum_{r:q_{ar}>0}q_{ar}\int_0^1
Q_{F_{ar}}(u)Q_{W_{ar}}(u)\,du.
\]
The sharp ATE interval for \(P_{\mathrm{rel}}\) has length \(D_H(g)\). Moreover, there are two probability laws \(P_-,P_+\in\mathfrak M(H,g,P_{\mathrm{rel}})\), as defined in \cref{obj:def:compatible-law-class}, such that
\[
\mathbb E_{P_-}[Y_1-Y_0]=\underline\mu_1-\overline\mu_0,
\qquad
\mathbb E_{P_+}[Y_1-Y_0]=\overline\mu_1-\underline\mu_0,
\qquad
\mathbb E_{P_+}[Y_1-Y_0]-\mathbb E_{P_-}[Y_1-Y_0]=D_H(g).
\]
\end{theoremv}

Assignment tilts the score distribution within a label differently in the two arms. Pairing the Bernoulli outcome with reciprocal assignment weights in opposite or matching quantile order produces the mean endpoints in \cref{obj:thm:all-label-ambiguity}. Its two compatible causal laws attain those endpoints under the same released outcome law, so the paired absolute-deviation expression gives the exact worst-case width. Because the expression depends on \(H\) and \(g\), it permits comparison of candidate releases without specifying an outcome regression. The result applies to arbitrary probability laws \(H\), including atomic laws, and to measurable releases with disconnected score cells.

For a concrete one-label release, take \(\varepsilon=1/4\), let \(H\) put mass \(1/2\) on each of the scores \(1/4\) and \(3/4\), and set \(g\equiv1\). Use the released law in \cref{obj:thm:all-label-ambiguity}, so outcomes are Bernoulli\((1/2)\) in each arm. In either arm the lower and upper attainable potential-outcome means are \(1/3\) and \(2/3\): assignment tilts the two scores, but the outcome can be matched to the larger or smaller reciprocal weight. The sharp ATE interval is therefore \([-1/3,1/3]\), with width \(2/3\). If each score instead receives its own label, the score is recovered from the label and the ambiguity width is zero. This example isolates uncertainty from within-label score variation before any trial sampling is introduced.

The zero-width case has a direct disclosure interpretation. The following characterization expresses it both as recovery of the score from the label and as concentration of the score within every positive-mass cell.

\begin{theoremv}{thm:universal-point-identification}[Universal point identification]
Let \(J\in\mathbb N\), let \(H\) be a probability law on the score interval \(\mathcal E=[\varepsilon,1-\varepsilon]\), and let \(g:\mathcal E\to\mathcal R_J\) be a measurable deterministic release into \(J\) labels. Under \cref{obj:ass:overlap}, write \(C_r=\{e\in\mathcal E:g(e)=r\}\) for each label \(r\). Each of the following is equivalent to \(D_H(g)=0\), where \(D_H(g)\) is defined in \cref{obj:def:worst-case-ambiguity}:
\begin{enumerate}
\item There exists a measurable map \(h:\mathcal R_J\to\mathcal E\) such that \(e=h(g(e))\) for \(H\)-almost every \(e\).
\item For every \(r\in\mathcal R_J\) with \(H(C_r)>0\), there exists \(x\in\mathcal E\) such that \(e=x\) for \(H\)-almost every \(e\in C_r\).
\end{enumerate}
\end{theoremv}

The absolute-deviation terms in \cref{obj:thm:all-label-ambiguity} vanish when each positive-mass cell carries a single score value. Score variation within a cell produces positive worst-case ambiguity, while score-degenerate cells yield universal point identification. For a finite-support score law with at most \(J\) positive-mass atoms, giving each atom a distinct label achieves point identification. With more atoms than labels, at least one positive-mass cell contains score variation, so universal point identification fails.
